\section{Part I: Paired Weight Codes}
\label{sec:part1}

\subsection{The hypothesis and the control}

The codec pairs two real weights as $z=w_a+iw_b$ and stores a three-bit
magnitude index and a five-bit phase index per pair -- the same eight index bits
two real 4-bit weights would take. An earlier draft of this work reported that
the code preserved frozen Qwen3.5-0.8B validation loss better than an
equal-byte real 4-bit control across all eighteen DeltaNet input projections,
with $\Delta$NLL $+0.009333$ against the control's $+0.021945$.

That comparison used a real 4-bit control whose codebook is a \emph{uniform}
16-level grid on $[-1,1]$. The polar code's effective codebook is not uniform: a
radius grid crossed with a phase grid produces annular sectors, so the polar
code was receiving a shape adaptation the control was denied. The asymmetry is
in the tuning, not in the representation.

We therefore add a matched control: the same scalar code with a per-tensor
16-level Lloyd--Max codebook fitted to that tensor's own normalized weight
distribution, and the identical per-row FP32 scale. The extra storage is 16 FP32
per tensor, under $10^{-4}$ bits per weight for these matrices.

\subsection{A ceiling for the whole family}
\label{sec:ceiling}

Before asking whether polar helps, we ask how much room exists at eight index
bits per pair, and whether the polar arrangement uses it well. We fit an
\emph{unconstrained} 256-point two-dimensional codebook to the pairs by
$k$-means. Any scheme that pairs two real weights and spends eight index bits on
the pair -- polar, Cartesian, rotated, or otherwise -- is a constrained special
case, so its distortion lower-bounds the entire family at this budget.

\begin{table}[ht]
\centering\small
\setlength{\tabcolsep}{4pt}
\caption{Weight reconstruction at eight index bits per pair, relative to the Lloyd-fitted
real-4 control (values below $1.000$ beat it). \texttt{vq2d} is an unconstrained 256-point
2-D codebook and lower-bounds distortion for every paired code at this budget.
Every entry uses the best of four pairings and a per-row FP32 scale.}
\label{tab:ceiling}
\begin{tabular}{llrrrrr}\toprule
Checkpoint & Tensor & real4-Lloyd MSE & real4-unif. & polar-legacy & polar-fitted & vq2d\\\midrule
Qwen3.5-0.8B & DeltaNet QKV L0 & 3.252e-06 & 1.308 & 1.194 & 1.065 & \textbf{0.759}\\
Qwen3.5-0.8B & DeltaNet QKV L8 & 3.135e-06 & 1.214 & 1.156 & 1.064 & \textbf{0.793}\\
Qwen3.5-0.8B & DeltaNet QKV L16 & 5.039e-06 & 1.240 & 1.165 & 1.059 & \textbf{0.795}\\
Qwen3.8-27B & DeltaNet QKV L0 & 2.563e-06 & 1.186 & 1.162 & 1.073 & \textbf{0.818}\\
Qwen3.8-27B & Attention Q L3 & 3.317e-06 & 1.217 & 1.154 & 1.060 & \textbf{0.791}\\
Qwen3.8-27B & MLP up L3 & 1.182e-06 & 1.169 & 1.163 & 1.081 & \textbf{0.817}\\
\bottomrule
\end{tabular}
\end{table}


Three facts follow from Table~\ref{tab:ceiling}, on every tensor measured across
two checkpoints.

\begin{enumerate}
\item \textbf{The reported advantage is an artifact of the control.} Polar beats
the uniform scalar code, as reported, but the Lloyd-fitted scalar code beats
polar on all six tensors.
\item \textbf{The polar family is structurally capped.} Even after we repair its
assignment rule, fit its radius grid, and sweep every magnitude/phase bit split
(Section~\ref{sec:polarfix}), the best polar code remains $6$--$8\%$ above the
fitted scalar code.
\item \textbf{The pairing idea is sound.} A free 2-D codebook sits $18$--$24\%$
\emph{below} the fitted scalar code at the same index budget. The headroom is
real; the polar parameterization fails to reach it.
\end{enumerate}

\subsection{Why repairing the polar assignment rule changes almost nothing}
\label{sec:polarfix}

The original quantizer rounds magnitude and phase independently, which is not
nearest-neighbour assignment in the $8\times32$ codebook. For a decoded phase
ray $\theta_p$ and a source $z=re^{i\theta}$,
\begin{equation}
 |z-g_m e^{i\theta_p}|^2 = r^2\sin^2\delta + (g_m - r\cos\delta)^2,
 \qquad \delta=\theta-\theta_p ,
\end{equation}
so the optimal radius level is the one nearest the \emph{projection}
$r\cos\delta$, not the one nearest $r$. We implemented exact nearest-neighbour
assignment over all 256 codepoints, verified in float64 against brute force
(maximum relative slack $2.9\times10^{-7}$, zero of 16{,}384 pairs suboptimal).
It reduced weight MSE by $0.035\%$.

The reason is quantitative and retrospectively explains a long series of null
results. With five phase bits the nearest-ray error satisfies
$|\delta|\le\pi/32$, hence $1-\cos\delta\le4.8\times10^{-3}$. The projection
correction is at most half a percent of the radius, while one step of the
three-bit radius grid is about $14\%$ of the row scale, so the correction almost
never changes which radius level wins.

Every phase-side refinement attempted across this project -- learned global
phase offsets, activation-weighted phase choice, output-error phase selection --
was operating on a quantity bounded to be small, which is why each produced
sub-noise, sign-unstable effects across layers and corpora. \textbf{The binding
constraint is the radius grid, not the phase grid.} The bit-split sweep confirms
it: $8\times32$ is optimal on every tensor, and moving bits toward the phase is
catastrophic ($4\times64$ roughly $3\times$ worse, $2\times128$ roughly
$14\times$ worse).

\subsection{Held-out loss, across three selection protocols}

Weight reconstruction does not settle the question; this project has repeatedly
found that weight-error rankings fail to predict task rankings. We therefore
replaced all eighteen projections simultaneously and evaluated the complete
WikiText-2 validation stream. Because the earlier draft selected each layer's
pairing by calibration NLL rather than by weighted error, and because that is a
stronger rule, we report both, plus a reproduction of the earlier draft's exact
calibration size.

\begin{table}[ht]
\centering\small
\caption{Frozen Qwen3.5-0.8B with all 18 DeltaNet input projections quantized
simultaneously, evaluated on all 2041 contiguous 128-target
WikiText-2 validation blocks. Every codec spends four index bits per real weight
and one FP32 row scale. Intervals are paired block bootstraps against BF16
(5{,}000 replicates). Lower is better.
$^{a}$the earlier draft's control; $^{b}$the matched control;
$^{c}$the earlier draft's method.}
\label{tab:main}
\begin{tabular}{lrrrl}\toprule
Method & NLL & Perplexity & $\Delta$NLL & 95\% CI\\\midrule
BF16 reference & 3.436670 & 31.083 & --- & ---\\
\midrule
real-4, uniform codebook$^{a}$ & 3.458469 & 31.768 & +0.021798 & $[+0.020841,\,+0.022752]$\\
real-4, Lloyd codebook$^{b}$ & 3.451102 & 31.535 & +0.014431 & $[+0.013626,\,+0.015235]$\\
\quad + sequential calibration & 3.451124 & 31.536 & +0.014454 & $[+0.013641,\,+0.015255]$\\
polar 3+5, decoupled rounding$^{c}$ & 3.456446 & 31.704 & +0.019776 & $[+0.018821,\,+0.020732]$\\
\quad + sequential calibration & 3.456446 & 31.704 & +0.019776 & $[+0.018821,\,+0.020732]$\\
polar 3+5, exact NN + fitted radii & 3.452258 & 31.572 & +0.015588 & $[+0.014736,\,+0.016421]$\\
\quad + sequential calibration & 3.451662 & 31.553 & +0.014992 & $[+0.014156,\,+0.015820]$\\
paired 2-D codebook (256 points) & 3.445063 & 31.345 & +0.008393 & $[+0.007720,\,+0.009069]$\\
\quad + sequential calibration & 3.445670 & 31.364 & +0.009000 & $[+0.008321,\,+0.009679]$\\
\bottomrule
\end{tabular}
\end{table}

\begin{table}[ht]
\centering\small
\caption{Paired head-to-head differences on the same validation blocks.
A negative value favours the first-named method.}
\label{tab:head}
\begin{tabular}{lrll}\toprule
Comparison & $\Delta$NLL & 95\% CI & Interval\\\midrule
polar (earlier) $-$ real-4 Lloyd & +0.005345 & $[+0.004282,\,+0.006407]$ & excludes 0\\
polar (repaired) $-$ real-4 Lloyd & +0.001156 & $[+0.000156,\,+0.002166]$ & excludes 0\\
paired 2-D $-$ real-4 Lloyd & -0.006038 & $[-0.006952,\,-0.005104]$ & excludes 0\\
paired 2-D $-$ polar (earlier) & -0.011383 & $[-0.012410,\,-0.010380]$ & excludes 0\\
\bottomrule
\end{tabular}
\end{table}


\begin{table}[ht]
\centering\small
\setlength{\tabcolsep}{4pt}
\caption{One protocol for every codec: per-layer candidate selection by calibration NLL
against an otherwise-BF16 model, 16 training windows, then cumulative application.
Codecs that can use activation statistics receive the identical BF16-measured statistic.
Complete WikiText-2 validation stream; paired block bootstraps against BF16.}
\label{tab:fair}
\begin{tabular}{lrrl}\toprule
Method & NLL & $\Delta$NLL & 95\% CI\\\midrule
real-4, uniform codebook (earlier draft's control) & 3.460834 & +0.024164 & $[+0.023088,\,+0.025239]$\\
polar 3+5 (earlier draft's method) & 3.452573 & +0.015902 & $[+0.015006,\,+0.016825]$\\
real-4, Lloyd codebook (matched control) & 3.445209 & +0.008539 & $[+0.007662,\,+0.009403]$\\
paired 2-D codebook & 3.445815 & +0.009145 & $[+0.008325,\,+0.009950]$\\
\textbf{paired 2-D codebook, activation-weighted} & 3.444363 & +0.007693 & $[+0.007029,\,+0.008398]$\\
\midrule
\multicolumn{4}{l}{\footnotesize Earlier draft's published polar figure: $+0.009333$.}\\
\bottomrule
\end{tabular}
\end{table}

\begin{table}[ht]
\centering\small
\setlength{\tabcolsep}{4pt}
\caption{Paired differences under the single protocol of Table~\ref{tab:fair}.
A negative value favours the first-named method.}
\label{tab:fairhead}
\begin{tabular}{lrll}\toprule
Comparison & $\Delta$NLL & 95\% CI & Interval\\\midrule
polar $-$ matched scalar control & +0.007364 & $[+0.006332,\,+0.008376]$ & excludes 0\\
polar $-$ paired 2-D (act.-weighted) & +0.008210 & $[+0.007205,\,+0.009186]$ & excludes 0\\
paired 2-D $-$ same, act.-weighted & +0.001452 & $[+0.000556,\,+0.002327]$ & excludes 0\\
matched scalar $-$ paired 2-D (act.-weighted) & +0.000846 & $[-0.000169,\,+0.001822]$ & \textbf{includes 0}\\
\bottomrule
\end{tabular}
\end{table}


The ranking is stable wherever the protocol can discriminate:

\begin{center}\small
\begin{tabular}{lrrl}\toprule
Protocol & polar & real-4 Lloyd & Discriminates?\\\midrule
Weighted-error selection, 16 windows & $+0.019776$ & $+0.014717$ & yes, Lloyd\\
NLL selection, 16 windows & $+0.015902$ & $+0.008539$ & yes, Lloyd\\
NLL selection, 4 windows, 3 seeds & $0.010945\pm0.000422$ & $0.010554\pm0.003349$ & \textbf{no}\\\bottomrule
\end{tabular}
\end{center}

The third row is the earlier draft's own calibration size, and we report that it
does \emph{not} decide the question. Over three seeds the two codecs are
statistically tied and on one seed polar wins outright. The matched control's
seed-to-seed spread is eight times larger, because its clip-factor search is a
finer decision than polar's four-way pairing choice and 512 calibration tokens
cannot resolve it. That is a narrow but genuine point in polar's favour: coarser
decisions are more stable when calibration data is scarce.

What the small-calibration sweep does establish is that all three of our polar
reproductions fall in $[+0.010515,+0.011358]$, so the earlier draft's published
$+0.009333$ lies outside the spread we could reproduce and is best read as a
fortunate draw from a high-variance selector.

\subsection{What survives: parity, not superiority}

Under a single protocol in which every codec receives per-layer NLL selection
and the codecs that can use activation statistics receive the identical
statistic, the activation-weighted 2-D codebook gives the best point estimate of
anything we tested. Two of its three comparisons are conclusive and one is not:

\begin{itemize}
\item It beats polar decisively, $+0.008210$ $[+0.007205,+0.009186]$.
\item Activation weighting is what makes the difference: the unweighted 2-D
 codebook is worse by $+0.001452$ $[+0.000556,+0.002327]$.
\item Against the matched scalar control the paired difference is $+0.000846$
 with interval $[-0.000169,+0.001822]$, which \emph{includes zero}.
\end{itemize}

We therefore claim parity, not superiority: the 2-D codec reaches the
well-tuned scalar code's regime and may exceed it, while the polar code is
significantly behind both. Replacing the polar grid with a learned codebook also
removes the obstacle that defeated our packed kernels -- decoding a polar index
requires a sine and a cosine per pair, which made decode-plus-GEMM
$3$--$4.5\times$ slower than resident BF16, whereas decoding a 2-D codebook
index is a 256-entry lookup that fits in shared memory. We have not written or
benchmarked that kernel, so this removes a known blocker rather than
demonstrating a speedup.

\subsection{Summary of Part I}

Polar magnitude--phase weight coding loses to a real 4-bit code once the scalar
code receives the same per-tensor codebook fitting, and it is structurally
capped several percent above that control regardless of how its assignment rule,
radius grid, and bit split are tuned. The measurement that establishes this also
locates the headroom: an unconstrained 2-D codebook over the same pairs sits far
below the scalar code at an identical index budget. The pairing idea was right
and the polar parameterization was wrong -- but the resulting codec reaches
parity with a well-tuned scalar code rather than beating it.
